\section{Yield Strategy Interface}
\label{sec:strategies}
% ============================================================

\subsection{Interface Definition}

Every yield source implements the \texttt{IYieldStrategy} interface:

\begin{lstlisting}[language=Solidity,caption=IYieldStrategy interface]
interface IYieldStrategy {
    function deposit(uint256 amount)
        external returns (uint256 shares);
    function withdraw(uint256 shares)
        external returns (uint256 assets);
    function totalAssets()
        external view returns (uint256);
    function currentAPY()
        external view returns (uint256);
    function harvest()
        external returns (uint256 harvested);
    function isActive()
        external view returns (bool);
    function riskTier()
        external view returns (uint8);
}
\end{lstlisting}

\begin{definition}[Strategy Properties]
A compliant strategy satisfies:
\begin{enumerate}[nosep]
  \item \textbf{Idempotent harvest}: Calling \texttt{harvest()} twice in succession yields zero on the second call.
  \item \textbf{Withdrawal guarantee}: \texttt{withdraw(s)} returns at least \texttt{convertToAssets(s) - slippageTolerance} assets.
  \item \textbf{Accounting accuracy}: $|\texttt{totalAssets()} - A_{\text{actual}}| \leq \epsilon$ where $\epsilon$ is bounded by oracle precision.
  \item \textbf{Risk classification}: \texttt{riskTier()} returns a value in $\{1, 2, 3\}$ corresponding to low, medium, or high risk.
\end{enumerate}
\end{definition}

\subsection{Strategy Taxonomy}\label{sec:taxonomy}

Table~\ref{tab:full_strategies} enumerates all 29 deployed strategies across 9 categories.

\begin{table*}[t]
\centering
\small
\caption{Complete yield strategy taxonomy: 29 strategies across 9 categories. APY ranges reflect Q1 2025 observed rates. Risk tiers: T1 = low (blue-chip, audited, \$1B+ TVL), T2 = medium (established, \$100M+ TVL), T3 = high (newer protocols, higher yield, elevated risk).}
\label{tab:full_strategies}
\begin{tabular}{@{}rllllrr@{}}
\toprule
\textbf{\#} & \textbf{Category} & \textbf{Protocol} & \textbf{Underlying} & \textbf{Mechanism} & \textbf{APY (\%)} & \textbf{Tier} \\
\midrule
\multicolumn{7}{l}{\textit{Liquid Staking}} \\
1 & Liquid Staking & Lido & ETH & Validator staking, stETH issuance & 3.2--3.8 & T1 \\
2 & Liquid Staking & Rocket Pool & ETH & Permissionless validators, rETH & 3.0--3.5 & T1 \\
\midrule
\multicolumn{7}{l}{\textit{Restaking}} \\
3 & Restaking & EigenLayer & ETH/stETH & AVS security, restaking rewards & 3.5--6.0 & T2 \\
4 & Restaking & Symbiotic & ETH/stETH & Flexible restaking, multi-asset & 3.0--5.5 & T2 \\
5 & Restaking & Karak & ETH/stETH & DSS restaking, risk isolation & 3.0--5.0 & T2 \\
\midrule
\multicolumn{7}{l}{\textit{Lending}} \\
6 & Lending & Aave V3 & ETH/USDC/USDT & Variable-rate lending pools & 2.0--8.0 & T1 \\
7 & Lending & Compound V3 & ETH/USDC & Isolated lending markets & 2.5--7.0 & T1 \\
8 & Lending & Morpho & ETH/USDC & Peer-to-peer rate optimization & 3.0--9.0 & T2 \\
9 & Lending & Euler V2 & Multi-asset & Modular lending, custom vaults & 3.5--10.0 & T2 \\
10 & Lending & Spark & DAI/ETH & MakerDAO-native lending & 4.0--8.0 & T1 \\
11 & Lending & Fluid & ETH/USDC & Unified liquidity layer & 3.0--12.0 & T2 \\
\midrule
\multicolumn{7}{l}{\textit{Stablecoin Yield}} \\
12 & Stablecoin & MakerDAO/Sky & DAI & DSR + Sky savings rate & 5.0--8.0 & T1 \\
13 & Stablecoin & Ethena & USDe & Basis trade (spot + short perp) & 8.0--15.0 & T3 \\
14 & Stablecoin & Frax & FRAX & AMO yield + staking & 4.0--7.0 & T2 \\
\midrule
\multicolumn{7}{l}{\textit{LP / DEX}} \\
15 & LP/DEX & Curve & Stablecoins & StableSwap AMM LP fees & 5.0--12.0 & T1 \\
16 & LP/DEX & Convex & CRV/CVX & Boosted Curve LP + CVX rewards & 8.0--18.0 & T2 \\
17 & LP/DEX & Pendle & PT/YT & Yield tokenization, fixed-rate & 5.0--20.0 & T2 \\
18 & LP/DEX & L2 DEX Aggregate & Multi-asset & Concentrated liquidity on L2s & 6.0--15.0 & T2 \\
\midrule
\multicolumn{7}{l}{\textit{Bitcoin Native}} \\
19 & BTC Native & Babylon & BTC & BTC staking for PoS security & 3.0--5.0 & T2 \\
20 & BTC Native & Lombard & BTC & LBTC liquid staking derivative & 3.5--6.0 & T2 \\
21 & BTC Native & SolvBTC & BTC & BTC yield aggregation vaults & 4.0--8.0 & T3 \\
22 & BTC Native & CoreDAO & BTC & Dual staking (BTC + CORE) & 3.0--12.0 & T2 \\
\midrule
\multicolumn{7}{l}{\textit{Solana}} \\
23 & Solana & Marinade & SOL & mSOL liquid staking & 5.0--7.0 & T1 \\
24 & Solana & Jito & SOL & JitoSOL + MEV rewards & 6.0--8.0 & T2 \\
25 & Solana & Kamino & SOL/USDC & Concentrated liquidity vaults & 5.0--10.0 & T2 \\
\midrule
\multicolumn{7}{l}{\textit{TON}} \\
26 & TON & Tonstakers & TON & tsTON liquid staking & 5.0--7.0 & T2 \\
27 & TON & Bemo & TON & stTON liquid staking & 5.0--7.0 & T2 \\
28 & TON & STON.fi & TON/USDT & DEX LP on TON & 8.0--15.0 & T3 \\
\midrule
\multicolumn{7}{l}{\textit{Perpetuals}} \\
29 & Perps & LPX Perps & Multi-asset & LP vault earns trading fees & 10.0--30.0 & T3 \\
\bottomrule
\end{tabular}
\end{table*}

\subsection{Yield Composition}\label{sec:composition}

Bridge token holders earn layered yield from multiple simultaneous sources. We formalize this as additive yield composition:

\begin{definition}[Yield Composition]
For a yield-bearing bridge token $yLA$, the total APY is:
\begin{equation}\label{eq:yield_composition}
  r_{\text{total}} = r_{\text{base}} + r_{\text{restaking}} + r_{\text{bridge}} + r_{\text{DeFi}}
\end{equation}
where:
\begin{itemize}[nosep]
  \item $r_{\text{base}}$: primary yield source (staking, lending, LP)
  \item $r_{\text{restaking}}$: additional yield from restaking derivatives
  \item $r_{\text{bridge}}$: share of Teleport bridge fees via xLUX
  \item $r_{\text{DeFi}}$: yield from using yLA as collateral/LP on destination chain
\end{itemize}
\end{definition}

\begin{remark}
The DeFi layer ($r_{\text{DeFi}}$) is not part of the vault's backing---it is additional yield earned by the token holder through their own DeFi participation on the destination chain. The vault only manages $r_{\text{base}} + r_{\text{restaking}}$, and $r_{\text{bridge}}$ flows through the xLUX staking mechanism.
\end{remark}

Table~\ref{tab:yield_composition} illustrates the yield stack for three representative bridge tokens.

\begin{table}[h]
\centering
\small
\caption{Yield composition for representative bridge tokens. Rates are Q1 2025 observed medians.}
\label{tab:yield_composition}
\begin{tabular}{@{}lrrr@{}}
\toprule
\textbf{Layer} & \textbf{yLETH} & \textbf{yLBTC} & \textbf{yLUSD} \\
\midrule
Base yield & 3.5\% & 4.5\% & 6.0\% \\
Restaking & 2.0\% & 2.0\% & --- \\
Bridge fees (xLUX) & 1.0\% & 1.0\% & 1.0\% \\
DeFi on Lux & 5--15\% & 5--15\% & 5--12\% \\
\midrule
\textbf{Total range} & \textbf{11.5--21.5\%} & \textbf{12.5--22.5\%} & \textbf{12.0--19.0\%} \\
\bottomrule
\end{tabular}
\end{table}

% ============================================================
